\summarychapter{7}{The School for Sinners}{The Power of Bishops}

\begin{inbrief}
In the apostolic age the Spirit moved freely among prophets, teachers and healers. By the middle of
the third century a single bishop led each church and held the power to forgive even the gravest sins.
The chapter traces the rise of the bishop's office and the crisis over sin after baptism that turned
the church from a society of saints into a school for sinners---the last element of catholic
Christianity.
\question{What if a Christian sins seriously after baptism? Are there sins for which the church has no
forgiveness, and who may forgive?}
\end{inbrief}

The first Christians knew a believer by one thing above all, rebirth through the Spirit, and baptism
was its sign. The apostles were the unchallenged leaders, yet prophets, teachers, miracle workers and
healers also acted in the Spirit's power, and Paul, even amid the rivalries at Corinth, would not
suppress these gifts. The refusal of Montanism, with its prophecies and its moral rigor, showed the
church as an institution more clearly than before. It had preached to the nations and shown its
universality; it had faced heretics and stated its orthodoxy; now it struggled with sin and developed
its episcopacy (from \emph{episcopos}, bishop).

\ssection{Leaders for the Churches}
Paul left local leaders in his churches: \emph{elders} or presbyters, also called bishops (overseers)
or pastors (shepherds), who taught converts, led worship and kept discipline, and \emph{deacons}, who
assisted them. So the apostolic age knew both traveling leaders moved by the Spirit and resident ones.
Soon after 100 Ignatius of Antioch writes as a matter of course of a single bishop in each church with
his presbyters and his deacons, through whose united ministry grace flows to the flock.
How this pattern arose is not known; probably one presbyter came to handle correspondence with other
churches, the funds for the poor, the defense of the faith against heretical teachers and the
Eucharist. It spread gradually---Alexandria had no single bishop until about 180---and unevenly:
small churches in Asia Minor and Africa each had a bishop, while in Gaul the bishop of a large town
supervised the surrounding congregations through presbyters.

By the late second century the bishop was the undisputed leader, and the struggle with the Gnostics
strengthened him. They claimed a secret wisdom handed down through a succession of teachers from
Jesus. Catholics answered by pointing to what the churches taught openly, to the rule of faith and to the bishops who presided where the apostles themselves had planted churches. Hegesippus, who journeyed from Palestine to Rome about the middle of the second century, heard the same teaching from every bishop and drew up succession lists for
Corinth and Rome; Irenaeus and Tertullian argued likewise that the succession of bishops guaranteed
the unbroken tradition of apostolic doctrine.

\ssection{What to Make of Bishops}
Nearly all Christians admit these changes took place; they disagree about their authority, and Shelley
sets out three answers. (1) The early leaders sinfully departed from a divinely given pattern, and the
changes should be reversed: the view of ``back-to-the-Bible'' and most reforming movements, which must
then decide how much of apostolic practice is permanent (elders, yes; the silence of women, too?).
(2) There was no divine pattern; the leaders used their liberty, and church government may change
with the times: the view of ``progressives,'' which risks a Christianity without absolutes. (3) The
Holy Spirit so guided the church that these developments are God's work and permanently binding: the
view of most Catholics, which must explain why the Spirit's guidance should stop with the early
centuries. These conflicting answers partly explain today's denominations.

As late as the 200s many believers suspected that the Spirit had left when the bishops took charge.
Earlier Christians had looked for the power of the Spirit in lives, not in an office. Paul described the
Spirit's work as the building up of the church and its fruit as love, joy, peace and self-control, and
the generations after him put moral demands first, tolerated no unholy members and expelled gross
sinners. The high morality of the congregations was a main argument of the apologists: Justin tried to
show that Christian goodness was a power actually exercised, and Athenagoras pointed to artisans and
old women who could not explain their doctrines but proved them by their deeds. Opponents agreed.
Pliny found nothing criminal in the Christians he examined, and the steadfastness of the martyrs
convinced Justin and won over onlookers. But early in the third century this moral fiber weakened.
Montanus had not been wholly wrong: by about 220 the churches and those who led them had changed from what they once were.

\ssection{What About Unforgivable Sins?}
For two centuries most Christians held that baptism canceled all earlier sins, and that three sins
committed afterward---sexual immorality, murder and denial of the faith (apostasy)---could be
forgiven by God but never by the church. The guilty were excluded from fellowship and from the Lord's
Supper, which Ignatius had called ``the medicine of immortality''; to be barred from it endangered
salvation. During the long calm of the early third century standards relaxed. The bishop of Rome,
Callistus (217--222), was the first to readmit penitents as a policy, including those guilty of
adultery and abortion. The church, he said, is like Noah's ark, with clean and unclean beasts; and the
bishop of Rome, being ``close to Peter,'' holds the keys to bind and loose---the first such claim by
a Roman bishop. Tertullian was horrified: ``We do not forgive apostates, and shall we forgive
adulterers?'' But the future lay with Callistus, and soon people asked: why not apostates too?

In 250 the emperor Decius, a general from the Danube who blamed the empire's troubles on Christian
``atheism,'' ordered every citizen to sacrifice to the Roman gods and obtain a certificate
(\emph{libellus}). It was the most violent persecution so far. Many complied or bought false
certificates; others went to prison or to death, and the bishops of Rome, Antioch and Jerusalem were among the victims. Decius wanted apostates rather than heroes, so Christians were tortured to make them say
``Caesar is Lord.'' Those who died were ``martyrs'' (witnesses), those who endured ``confessors,''
those who gave way the ``lapsed.'' The storm ended in 251 when Decius fell in battle against the
Goths. Then the question of readmission became urgent, for in places three-quarters of a congregation
had lapsed. Cyprian, bishop of Carthage, stated what was at stake: ``Outside the church there is no
salvation.'' Yet was not apostasy the unforgivable sin against the Holy Spirit?

\ssection{The Birth of Saints}
Martyrs and confessors were held in awe. Martyrdom, the baptism of blood, was a Christian's highest
glory; the martyrs' names were recorded and their ``birthdays'' into eternal life kept each year at
their tombs. The saints were on their way. In Carthage some held that the confessors, by their
courage, had received power from the Spirit to absolve sins and could cover the faults of the lapsed
with their own merits, and they pressed Cyprian for a general pardon. He refused and proposed
readmission by degrees of guilt: leniency for those who sacrificed only under extreme torture, the
severest penalty for those who went willingly. The church accepted this and created a graded system of
penance. After a period of sorrow, the penitent came before the congregation in sackcloth and ashes,
and the bishop laid hands on him as a sign of restoration. The confessors' idea did not die: it came back centuries later in Rome's teaching about a treasury of merit and in indulgences, which likewise credit the merits of the saints to sinners in need.

The strict tradition found its voice in Rome. The presbyter and theologian Novatian argued that the
church cannot forgive murder, adultery or apostasy and can only pray for God's mercy at the last
judgment. Another presbyter, Cornelius, held that the bishop can forgive even grave sins. Novatian
defended the older idea of the church as a society of saints; Cornelius stood for the new idea of the
church as a school for sinners. The majority elected Cornelius bishop; Novatian's minority formed a
network of small congregations that regarded the catholic churches as polluted---perhaps rightly,
Shelley remarks, since these now offered unlimited forgiveness. From then on catholics had, beside
baptism, a second sacrament, still without fixed form: penance. With the bishops' power to forgive
sins, catholic Christianity was complete.

\begin{keyfacts}
\keypeople{Ignatius of Antioch; Hegesippus; Irenaeus; Tertullian; Justin; Athenagoras; Callistus;
Decius; Cyprian of Carthage; Novatian; Cornelius.}
\keydates{%
  \kd{c. 180}{Alexandria receives a single bishop}
  \kd{217--222}{Callistus, bishop of Rome, readmits penitent adulterers and claims Peter's keys}
  \kd{c. 220}{Decline of moral discipline in the churches evident}
  \kd{250}{Decius orders universal sacrifice; many Christians lapse}
  \kd{251}{Decius killed by the Goths; dispute over readmitting the lapsed}}
\keyterms{Episcopacy; presbyter (elder), bishop, deacon; apostolic succession; rule of faith; apostasy;
\emph{libelli}; martyr, confessor, lapsed; penance; treasury of merit; Novatianists.}
\end{keyfacts}
